\section{Conclusion}\label{sec:conclusion}

The affine residue permutation and the interval identity retain different
parts of the finite fold.  The first resolves each fiber and yields the
partition-difference formula; the second identifies the complete multiset
with two consecutive Fibonacci layers and exposes the generator orbits.
Together they pass from finite arithmetic identities to local stabilization
and the critical renewal law without losing the window endpoints.
The same interval coordinates, combined with factorization at even
Zeckendorf zero-runs, also isolate the complete second-largest fiber level:
its value satisfies the extremal two-step recurrence, and its six- or
eight-fold location set is obtained from explicit concatenations of
maximizing source blocks.

At exponential scale, the freezing corner leaves a whole interval of
nonexposed slopes.  Convex duality determines the candidate affine rate but
does not prove its lower bound.  The orbit-filling map supplies that bound by
moving exposed source generators into every larger target layer while
preserving multiplicity, completing the finite-window large deviation
principle across the corner.
